% LADDER — JACoW-formatted conference paper (IPAC-style contributed paper).
% Class: jacow.cls v3.01 (vendored alongside; latest from
% https://github.com/JACoW-org/JACoW_Templates).
% Compile: pdflatex (twice) or drop both files into Overleaf.
% Page option: choose a4paper/letterpaper per the host conference.
% TODO before submission: author full name/initials, co-authors,
% funding footnote text per LBNL guidance, and LBNL/DOE release approval.
\documentclass[letterpaper]{jacow}

\usepackage[utf8]{inputenc}
\usepackage[T1]{fontenc}
\usepackage{graphicx}
\usepackage{booktabs}
\usepackage{url}

\begin{document}

\title{LADDER: HUMAN-FIRST, VENDOR-NEUTRAL GENERATION OF VERIFIED PLC PROGRAMS}

\author{N. Saqib\thanks{nusaqib@lbl.gov}, Lawrence Berkeley National Laboratory, Berkeley, CA, USA}

\maketitle

\begin{abstract}
Accelerator facilities depend on programmable logic controllers (PLCs)
whose programs live inside proprietary, binary, vendor-specific project
files: hard to review, hard to diff, and locked to one toolchain.
LADDER is an open toolkit that inverts this: control logic is authored
once as a declarative, schema-validated intermediate representation (IR)
of domain elements --- interlocks, redundant (1oo2) inputs, search/arming
chains, annunciator groups, timers, state machines, PID loops --- and
everything vendor-specific is generated deterministically. One IR
document drives five backends (Siemens TIA Portal via the Openness API,
Rockwell L5X, PLCopen XML, Beckhoff TwinCAT, and plain IEC~61131-3
text), a scan-accurate simulator with declarative acceptance scenarios,
a documentation package, and a formal-verification path that emits SMV
models with safety theorems derived automatically from each element's
semantics and proves them with nuXmv. Projects are self-contained git
repositories with pinned toolchains and continuous integration.
Authoring is organized as a human-gated loop over this deterministic,
LLM-free core: an assistant (any model) interviews the engineer for the
ground truth only they have and drafts the design map, IR, and
scenarios; the machine judges every draft with the same gates a human
edit would face; the engineer reviews and signs off. We describe the IR, the lowering and
verification pipeline, and a case study in which a personnel protection
system's fail-safe PLC program was reconstructed exactly --- including
certified F-instruction instances and the operator HMI --- from
version-controlled design sources with a single command.
\end{abstract}

\section{Introduction}

PLC programs at accelerator facilities concentrate three risks. First,
\emph{opacity}: the source of truth is a binary project file that
cannot be meaningfully diffed or reviewed in a pull request. Second,
\emph{vendor lock-in}: logic expressed as one vendor's ladder diagrams
cannot move to another platform, and often not even across versions of
the same tool. Third, \emph{informal change control}: modifications are
made interactively in an engineering IDE, so the actual safety-relevant
logic and its documentation drift apart.

Laboratories have long mitigated the symptoms by generating PLC code
from structured inputs~\cite{unicos,plcfactory,pyplc}, and CERN's
PLCverif has shown that model checking real PLC programs is practical
in production, including on a personnel safety system~\cite{plcverif15,
plcverif-status}. More recently, large language models (LLMs) have been
applied to Structured Text generation with compiler and model-checker
feedback~\cite{llm4plc,agents4plc}. LADDER combines these threads
differently: a vendor-neutral IR with domain semantics is the single
reviewed artifact, and generation, simulation, formal proof,
documentation, and multi-vendor deployment are all derived from it
deterministically. Authoring inverts the LLM-first framing: the
assistant's output is constrained, schema-validated YAML --- never
vendor code --- judged by the same deterministic gates as a human
edit, with the human supplying ground truth (signal senses, safety
philosophy, acceptance stories) and deciding at review gates. The
pipeline also runs entirely without a model.

\section{Related Work}

Config-to-code generation is established lab practice: UNICOS-CPC
instantiates a standard object model into Siemens and Schneider
code~\cite{unicos}; ESS's PLC Factory and its successor generate
Siemens blocks and EPICS interfaces from device
databases~\cite{plcfactory}; ALBA correlates equipment databases with
protection-logic templates~\cite{pyplc}. On the verification side,
PLCverif builds scan-accurate models from existing Siemens code and
checks requirements with nuXmv and other backends~\cite{plcverif15,
plcverif-status}; Arcade.PLC~\cite{arcade} and recent BMC-based
tools~\cite{esbmcplc} verify existing code similarly. The closest
single system in spirit is PLCspecif, a research-prototype formal
specification language generating Siemens code and supporting
verification~\cite{plcspecif}. LLM4PLC~\cite{llm4plc} and
Agents4PLC~\cite{agents4plc} generate free-form Structured Text inside
feedback loops. Vendors now endorse text-based, CI-built PLC workflows
(Siemens SIMATIC AX; Rockwell's CI/CD reference architecture with the
Logix Designer SDK~\cite{at002}).

To our knowledge, no prior system combines a vendor-neutral
\emph{semantic} IR, deterministic lowering to multiple real vendor
toolchains, and safety theorems \emph{generated automatically from
element semantics} rather than written by hand; that combination is
LADDER's contribution.

\section{The Intermediate Representation}

A LADDER project is YAML: user-defined types, a typed tag list with
explicit directions, and programs whose logic is a sequence of
\emph{elements}. Elements carry domain semantics, not syntax:

\begin{itemize}
\item \texttt{interlock}: permissive chain with latching trip and a
  named manual reset; the output is a permit.
\item \texttt{dual\_channel}: 1oo2 evaluation of a redundant input pair
  with discrepancy-time monitoring, a latched fault, and an
  acknowledge contract mirroring PROFIsafe F-DI behavior.
\item \texttt{search\_chain}: a search/arming sequence in which
  stations latch on the rising edge of their key, in walk order, only
  while a precondition holds; any breach clears the whole chain in one
  scan, and the element deliberately has no input that can restore a
  completed search.
\item \texttt{alarm\_group}: an ISA-18.2-style annunciator with
  per-member latches, shared horn and acknowledge, and first-out
  capture.
\item \texttt{alarm}, \texttt{timer}, \texttt{state\_machine},
  \texttt{scale}, \texttt{pid} (clamping anti-windup, bumpless
  enable), \texttt{assign}, plus expandable equipment patterns
  (e.g.\ motor starter) and a raw escape hatch.
\end{itemize}

Fail-safe conventions are enforced by validation and lint rules
(machine-actionable issue codes): Boolean inputs mean
1\,=\,OK/healthy, interlock outputs are permits, every latching
element names its reset or acknowledge, and hardware addresses are
banned from the IR --- they live in per-vendor IO maps. Programs may
declare any of the five IEC~61131-3 languages as their
\emph{presentation}; an expressibility check rejects logic that does
not fit the chosen language. Larger systems split the IR into a
reviewable directory (types, tags, one program per file, filename
order defining scan order).

\section{Deterministic Lowering and Backends}

Validation is followed by lowering to a small neutral statement AST
(assignments, conditionals, timer calls) with synthesized state
variables for edges, latches, and timers. Backends render this AST:
Siemens (SCL sources plus a headless build script that drives TIA
Portal Openness to an openable, compiled project), Rockwell (a
complete L5X with native ladder rung text and TIMER tags), PLCopen
tc6 XML (LD/FBD/SFC graphic bodies), Beckhoff TwinCAT source files,
and plain IEC~61131-3 text. Backends are versioned
(\texttt{siemens@21}, \texttt{rockwell@36}) so tool-generation quirks
stay behind one switch, and third-party backends load through a
plugin entry point.

Two vendor-free proofs run in continuous integration on every commit:
the emitted IEC text (ST, IL, and textual SFC) is compiled by the open
matiec compiler, and the emitted PLCopen XML is validated against the
official tc6 XSD. Ladder rendering uses a neutral rung model with
order-safety guards, so a rung realization that would change scan
semantics is rejected rather than silently emitted.

\section{Verification}

Three layers of evidence are generated from the same IR. First, a
\emph{scan-accurate simulator} executes the lowered AST with real
timer semantics; declarative scenario suites (set, pulse, run,
expect) are the project's acceptance gate and run in CI, with JUnit
output. Second, \emph{static lint} flags design smells such as
unreset latches or unread inputs. Third, \texttt{ladder model} emits
an SMV model per program via symbolic execution --- timers soundly
over-approximated so proofs hold for every preset --- and
\emph{auto-generates safety theorems from each element's semantics}:
an interlock output implies all its permissives; a dual-channel
output implies both channels; no acknowledge path can set a search
latch; a completed search implies every station latched in order.
User-supplied invariants join as
\texttt{given}/\texttt{always} pairs. nuXmv~\cite{nuxmv} proves the
result with BDDs, or IC3 for large programs; in the case study below,
79 theorems are proved, the largest models ($\sim$250 state bits)
via IC3.

\section{Self-Contained Projects}

Plant logic lives outside the toolkit, one git repository per machine,
scaffolded with a working starter. Each project vendors the exact
LADDER toolchain as a git submodule installed into a project-local
environment by a bootstrap script, and the manifest's version gate
refuses to run under any other toolchain --- a fresh machine needs
only git and Python to reproduce every artifact. \texttt{ladder docs}
generates the documentation package (requirements with one SHALL
statement per element, software specification, conventions, developer
and operator manuals, and a verification report) from the same IR, so
documentation cannot drift from logic. The LLM harness is
model-agnostic: \texttt{ladder prompt --intake} turns any chat model
into a design-intake interviewer that asks the engineer for exactly
the ground truth the machine cannot derive; \texttt{ladder generate}
closes the draft--validate--simulate loop mechanically around any
model CLI; and a six-task scenario-scored benchmark compares models.
Nothing in the pipeline requires a model.

\section{Case Study: A Personnel Protection System Reproduction}

As a forcing function against a real safety program, an operating
storage-ring personnel protection system's PLC software (Siemens
fail-safe, TIA Portal) was reverse-engineered into design sources ---
CSV hardware and signal sheets plus LADDER IR --- and reconstructed
exactly on a newer tool version (V19 reference, V21 rebuild). One
command rebuilds, from nothing: an F-CPU with 16 remote-IO stations
and 76 PROFIsafe modules; the as-built fail-safe data-block and type
tree; 144 tags; three generated F-LAD blocks containing 66 certified
safety-instruction instances (emergency stop, safety door, 1oo2
evaluation) matching the reference count exactly; and a 14-screen
operator HMI. The generated rungs are proved against a behavioral
simulator before import, the project compiles with zero errors, and
the portable path proves 79 safety theorems over the same logic. Two
Openness API limitations (PROFIsafe destination addresses and
integrated HMI connections are GUI-only) are documented as manual
steps --- matching decisions recorded in the original project. The
reproduction lives in a private repository; it is a study of
reproducibility, \emph{not} a certified or deployable safety system.

\section{Status and Outlook}

LADDER is open-source-ready (BSD-style license with a safety
disclaimer, contribution and security policies, CI); the IR is at
version 0.2 with a semver/RFC process toward a frozen 1.0. Planned
work includes requirement-pattern properties and counterexample
replay as failing scenarios (following PLCverif's
practice~\cite{plcverif-status}), mutation testing of scenario and
theorem strength, EPICS interface artifacts from the same IR, and
Rockwell/TwinCAT reverse adoption. Generated logic must always be
reviewed by qualified controls engineers; LADDER is engineering
automation and evidence generation, not certification.

\section{Acknowledgments}

This work was supported by the U.S. Department of Energy under
Contract No.\ DE-AC02-05CH11231.
% TODO: confirm funding/footnote wording per LBNL publication guidance.

\begin{thebibliography}{9}

\bibitem{unicos}
E. Blanco Vi\~nuela \emph{et al.}, ``UNICOS CPC6: automated code
generation for process control applications'', in \emph{Proc.
ICALEPCS'11}, Grenoble, France, 2011.

\bibitem{plcfactory}
G. Ulm \emph{et al.}, ``PLC factory: automating routine tasks in
large-scale PLC software development'', in \emph{Proc. ICALEPCS'17},
Barcelona, Spain, 2017.

\bibitem{pyplc}
S. Rubio-Manrique \emph{et al.}, ``Integration of PLCs in Tango
control systems using PyPLC'', in \emph{Proc. ICALEPCS'15},
Melbourne, Australia, 2015.

\bibitem{plcverif15}
D. Darvas, B. Fern\'andez Adiego, and E. Blanco Vi\~nuela,
``PLCverif: a tool to verify PLC programs based on model checking
techniques'', in \emph{Proc. ICALEPCS'15}, Melbourne, Australia,
2015, paper WEPGF092.

\bibitem{plcverif-status}
I. D. L\'opez-Miguel \emph{et al.}, ``PLCverif: status of a formal
verification tool for programmable logic controllers'', in
\emph{Proc. ICALEPCS'21}, Shanghai, China, 2022.

\bibitem{llm4plc}
M. Fakih \emph{et al.}, ``LLM4PLC: harnessing large language models
for verifiable programming of PLCs in industrial control systems'',
in \emph{Proc. ICSE-SEIP'24}, Lisbon, Portugal, 2024.

\bibitem{agents4plc}
Z. Liu \emph{et al.}, ``Agents4PLC: automating closed-loop PLC code
generation and verification in industrial control systems using
LLM-based agents'', arXiv:2410.14209, 2024.

\bibitem{plcspecif}
D. Darvas, E. Blanco Vi\~nuela, and I. Majzik, ``PLC code generation
based on a formal specification language'', in \emph{Proc. IEEE
INDIN'16}, Poitiers, France, 2016.

\bibitem{arcade}
S. Biallas, J. Brauer, and S. Kowalewski, ``Arcade.PLC: a
verification platform for programmable logic controllers'', in
\emph{Proc. ASE'12}, Essen, Germany, 2012.

\bibitem{esbmcplc}
K. Sun \emph{et al.}, ``ESBMC-PLC+: a unified IEC 61131-3 formal
verification framework'', arXiv:2606.23870, 2026.

\bibitem{at002}
Rockwell Automation, ``Industrial DevOps: CI/CD for Logix control
systems'', Application Technique LOGIX-AT002, 2023.

\bibitem{nuxmv}
R. Cavada \emph{et al.}, ``The nuXmv symbolic model checker'', in
\emph{Proc. CAV'14}, Vienna, Austria, 2014.

\end{thebibliography}

\end{document}
